\chapter{Conclusion}\label{chapter:conclusion}

This chapter closes the thesis by synthesizing the retrieval and patching results from Chapter~\ref{chapter:results}, stating what they suggest, and outlining suggestions for future work.

\section{Retrieval}

No embedder dominates retrieval performance across the two datasets tested. 
The Semantic (CodeBERT) embedder is the best on the AutoPatch dataset (0.877 hit@1, Table~\ref{tab:pca_l2}), with the unsupervised Topological Structure embedder close behind (0.822) despite requiring no pretrained language model and no GPU.
Breaking down the performance by CWE, Semantic embedding consistently performs well across categories, demonstrating better generalization, whereas Topological Structure-based embeddings exhibit heterogeneous performance.
On CVEFixes, conversely, the hybrid Semantic+Structure retriever leads with 65.17\% CWE accuracy and 53.93\% CVE accuracy (Section~\ref{sec:cvefixes-retrieval}), while Topological Structure shows 23.60\% CWE accuracy.
Given the difference between the datasets, it is possible to infer the sharpest capabilities of each embedder. 
The Structural embeddings excel when the vulnerability signal is a local, execution-order or control/data-flow pattern, as seen with concurrency bugs (100.0\% CVE/CWE hit rate on AutoPatch's Concurrency category, Table~\ref{tab:autopatch_by_category_scaled}) and arithmetic-into-allocation pipelines (CWE-190), the latter identified correctly by WL/GIN kernels even with frozen, shallow weights.
Semantic embeddings perform best when the signal resides in the domain vocabulary or is non-local, missing synchronization across functions or configuration-dependent triggers.

Regarding the pipeline, we identified that $L_2$-normalization is a major preprocessing component; its removal degrades hit@1 to 0.136 for Topological Structure, 0.055 for GIN, and 0.000 for WL (Table~\ref{tab:pca_l2}). PCA was characterized as a second-order effect component on the performance impact. 

The composition of the dataset also shows great importance. The CWE-362 subset, for example, with the Semantic+Structure pipeline, is brittle. For this embedding, the structural feature vector is built from regex matches against standard locking primitives (e.g., \texttt{mutex\_lock}, \texttt{spin\_lock}); when a codebase uses a custom primitive (\texttt{ip\_setup\_cork()}), that feature group goes to zero, and the vector becomes near-constant noise. Concatenating the structural vector with the CodeBERT embedding and reducing dimensionality with PCA dilutes the semantic signal. 
A hand-crafted structural detector is only as good as its pattern list.
On the one hand, this strategy works well for codebases standardized by the programming industry or internal directives; however, it is difficult to adapt and apply at scale. 
Additionally, once the patterns are identified, this vector can be enriched with CPG findings, making it a more efficient and highly explainable approach for vulnerability classification.

As the variant analysis results showed, retrieval success varies with the code structure. This was seen with the hand-augmented dataset, which achieves much lower performance than the LLM-augmented dataset and therefore retains the main structure. The practical consequence is that classification performance may vary with code style and quality, hindering the adoption of such strategies.

The main finding of retrieval experiments is the competitiveness of unsupervised Topological Structure to classify vulnerabilities. This is an important option for legacy or domain-specific languages where pretrained Language models, such as a CodeBERT-class model, have no coverage, while it remains viable (0.822 hit@1).
Besides not needing training, these embedder models can run on CPU, allowing a lightweight and cheaper approach to work early in the pipeline as a filter or as a fallback for a semantic retriever.

Another finding is that not all CWE classes suit the retrieval pipeline. 
Categories like Concurrency 
% (perfect topological retrieval and the strongest single lift, from 26.67\% to 73.33\% hit@1, for CWE-787's Semantic-vs-Structure+Semantic comparison in the CVEFixes retrieval chapter) 
and CWE-190 arithmetic-allocation bugs most clearly benefit from structural context over Semantic-only, whereas Memory Safety and Resource Management favor Semantic over pure structure
Conversely, NULL Pointer Dereference (CWE-476) is a case where the Semantic model's wide IQR of false positives is a known, documented weakness. This could be explained by the fact that a missing null check is lexically and sequentially near-identical to code where no check was ever needed, so no embedder investigated here reliably separates the two.

 %TODO discuss the 787 failure issue 
% CWE-787's near-total failure across every model and pipeline warrants a dedicated study of what specifically breaks -- bounds-check synthesis, buffer-size reasoning, or something dataset-specific to how CWE-787 examples were sourced -- since "everything fails at this uniformly" is itself informative but not yet explanatory.
 
\section{Patch Generation}

The pipeline experiment results also demonstrate that retrieval quality and patch quality are independent, debunking a hypothesis held at the beginning of this work.
Whereas the retrieval results showed that some approaches are better suited to a given vulnerability type, there is no clear "model specialist per CWE" in the patching quality results. 
Instead, we found a model-dependent interaction between retrieval and the LLM as a judge of patch validity, an unexpected finding in the patching chapter.
Retrieval helps GPT-5.3-Codex modestly pumping the baseline 21.3\% valid to Semantic 42.5\%, (see Table~\ref{tab:llmjudge_codex53}), helps DeepSeek-V4-Flash substantially pushing baseline 21.7\% to Semantic+Structure $\sim$27\%, as shown in Table~\ref{tab:deepseek_judge_summary}.
However, retrieval augmentation hurts GPT-4o, reducing the Tool Calling without retrieval from 20 valid patches to 5 on average for Semantic and Semantic+Structure retrieval (see Table~\ref{tab:llmjudge_gpt4o}), performing worse than the prompt-only baseline. 
This could be explained by the increase in the input length and complexity, since a code example is added to the prompt. Besides, the retrieval errors may not be as easily recovered as with the programming and security specialist model GPT-5.3-Codex.
This finding suggests that LLM choice is more complex than buying the more expensive model, as it simultaneously reinforces the latest indications of specialists' models usage. %TODO: cite works that say that
 
A qualitative analysis of the GPT-5.3-Codex experiments found that, for a query in which the retriever returned the wrong CVE, the LLM generated a technically coherent patch for a different, correctly diagnosed vulnerability. Still, its own chain of thought explicitly claimed to address the CWE stated in the (wrong) retrieved prompt. 
Even though the patch was correct, the dissimulation raises questions about the model's trustworthiness in production and the need to increase harnesses or, ultimately, specialists, to ensure the system's safety, as it is unclear in which direction and to what extent this deceitful behavior can occur.

Another research concern faced in this thesis, which also affects engineering teams, is the silent failure mode that may pass undetected by aggregated metrics, even when measurements capture a different spectrum of the evaluated object. 
As discussed in the Chapter \ref{chapter:results} on the baseline, empty fixes may result in inflating lexical scores that calculate F1 based on code addition and removal. For example, Codex's near-zero \texttt{hunk\_added\_f1} (0.011) against a high \texttt{hunk\_removed\_f1} (0.719) indicates that the model deletes vulnerable code without replacing it, resulting in an empty fix. 
Consequently, aggregate lexical metrics (BLEU-4, BERTScore, ROUGE-L, even whole-file Hunk F1) can look strong even when the patch is functionally empty, incorrect, or accompanied by a justification that does not match what the model actually did. 
In these cases, having a dataset with detailed annotations, e.g., indicating what should be added or removed, could improve handling of failure cases. Still, such labeling is expensive and hard to find for our use case.
Therefore, LLM-Judge validity was treated as the decisive metric throughout, despite its own limitations.
 
The adoption of a RAG pipeline and its configuration depend on the model adopted. 
For GPT-5.3-Codex and DeepSeek-V4-Flash, retrieval (particularly Semantic for Codex) produced a real gain in validity over tool-calling without retrieval and is worth the added cost. 
For GPT-4o, retrieval degraded validity below even the prompt-only baseline in the CVEFixes setting. For that model, tool-calling without retrieval is the better-performing and cheaper option.

Another important factor to consider when choosing a model is its error rate. 
For Codex, almost all generated patches were well-formed, and the reductions in validity rate were due to semantic errors. 
On the complex CVEFixes dataset, DeepSeek's patch generation validity rate is low (12--35 of 89 records produced a non-empty patch for the retrieval pipelines; see Table~\ref{tab:deepseek_judge_summary}) due to output-formatting failures. 
Using this model in production might require increased engineering effort to increase the patch generation rate, such as creating new tools, exploring new pipeline architectures, and implementing error-recovery mechanisms.

One approach to further reduce failure modes is to compile or run the generated patch against a test harness or a proof-of-concept exploit. 
As of today, the pipeline only performs text-similarity evaluation, and validity is determined by an LLM reading the diff. A sandboxed build-and-test step, e.g., compile success, existing unit tests still pass, PoC no longer triggers the CVE where one exists, would convert the LLM judgment reliability and could plausibly also serve as an agent feedback signal during generation rather than only a post-hoc filter. This would concomitantly address a real limitation of the evaluation methodology itself.
 
Augmenting context by grounding the model directly in the Code Property Graph, rather than only in linguistic prompt context, could help avoid the seen hallucinated justifications. 
% TODO: rephrase this
As a model that is shown a retrieved CVE only as text can misattribute its own reasoning to that text, whereas a model whose retrieved evidence is a structured graph object or whose patch is checked against one after the fact has a harder time producing a plausible but false causal reasoning, because the check does not rely on the model's own narration.

Regarding the cost, GPT-5.3-Codex is a cost vs. performance leader among the two proprietary models: on CVEFixes, it costs \$0.078--0.099 per entry with retrieval versus GPT-4o's \$0.244--0.262 (Table~\ref{tab:cvefixes_generation_cost_summary}), a $\sim$2.5$\times$ gap, while also producing higher BERTScore and far higher validity.
On AutoPatch, the gap narrows to $\sim$2.1--2.2$\times$, but Codex remains the winner (Table~\ref{tab:generation_cost_summary_condensed}).
For the open-weights model, DeepSeek-V4-Flash is 8--17$\times$ cheaper than Codex and GPT-4o, respectively, but this cost is amortized by the high error rate, meaning it does not account for patching the entire test set, since many entries had their pipelines fail.
 On the single-function AutoPatch dataset, the formatting bottleneck largely disappears (74-93\% completion), suggesting that DeepSeek's reliability is dataset-complexity-dependent, whereas Codex's is not.
Summing up, retrieval augmentation is not cost-neutral, and its cost is not always justified by the outcome. 
For Codex and DeepSeek, paying more for retrieval improves validity. For GPT-4o, retrieval is simultaneously the most expensive configuration (\$0.220-0.260/entry on CVEFixes) and the worst-performing one (4-6 valid patches, versus 20 for tool-calling without retrieval).
For future implementations, as context augmentation increases costs, this must be justified for each backbone model relative to its no-retrieval baseline.
Besides, a cheaper model with a higher completion rate is not automatically the better production choice if a meaningful share of its "successes" are unparseable or format-broken outputs rather than absent generation.

\section{Takeaways}\label{sec:conclusion_takeaways}

First, the embedding architecture is an auxiliary component as no embedder wins unanimously and rankings flip between datasets. 
It is important to reinforce that the AutoPatch numbers most favorable to Topological Structure are the ones most likely to be inflated by dataset construction. 
Second, base model capability dominates retrieval or prompting choices for patch validity, which means choosing a stronger backbone outweighs investment in retrieval sophistication. 
Third, retrieval augmentation is not a safe default; it is a per-model decision that can help (Codex, DeepSeek) or actively hurt (GPT-4o), and should not be treated as a universal upgrade according to the obtained results.
 
 
\section{Future Work}
 
The concrete next steps, in rough priority order: (1) a sandboxed compile-and-test verification stage, replacing or checked against LLM-Judge validity, ideally reused as an agent feedback signal during generation, not only as a post-hoc filter; (2) a Generator/Reviewer dual-agent architecture in which the Reviewer checks the patch's CPG against known-secure topologies rather than the Generator's self-reported reasoning, targeting the hallucinated-justification failure mode directly; (3) supervised fine-tuning of the GIN embedder to test whether a trained structural model can close or reverse the gap with CodeBERT-based retrieval on CPU-only infrastructure; (4) a controlled ablation isolating how much of AutoPatch's retrieval advantage is a dataset-construction artifact of topology-preserving augmentation; (5) replacing the regex-based structural pattern detectors (synchronization primitives being the clearest failure case) with a learned or more comprehensive alternative; and (6) extending the model-dependence finding on retrieval's helpfulness to additional backbones, since right now it rests on three models and one plausible but unconfirmed explanation.

The C-to-C++ degradation measured with Hunk F1 dropping needs its own investigation into diff/hunk fragmentation under RAII and template-heavy code. Finally, the model-dependent retrieval performance, it should be replicated on at least one more backbone to be able to judge its generalization or if it is a GPT-4o-specific behavior.

% Extend the pipeline to vulnerability identification, not only classification. As it could give a complementary signal to the LLM patching loop and avoid regression.  
